% Hodina 67, rozpis řádek 67. Zdroje: učebnice
% ucebnice/02_Funkce_a_rovnice/Linearni_funkce_a_rovnice/Lineární funkce.tex,
% motivace rovnoměrnou závislostí a aplikační úlohy; lokální Realisticky.cz
% 020115_Slovni_ulohy_na_linearni_funkce/020115_Lekce.pdf – metodická reference.
% Úlohy a graf jsou vlastní tvorba.
\documentclass[11pt,a4paper]{article}
% Spolecny styl pro ucitelske pripravy MAT (Work/MAT)
% Verze 1.1, 2026-10-23
\NeedsTeXFormat{LaTeX2e}
\ProvidesPackage{priprava}[2026/10/23 v1.1 Sablona ucitelskych priprav MAT]

% Zakladni typografie a cestina
\RequirePackage[T1]{fontenc}
\RequirePackage[utf8]{inputenc}
\RequirePackage[czech]{babel}
\RequirePackage{lmodern}
\RequirePackage{microtype}
\RequirePackage{amsmath,amssymb}
\RequirePackage{xcolor}
\RequirePackage{enumitem}
\RequirePackage[most]{tcolorbox}
\RequirePackage{tabularx,booktabs,array}
\RequirePackage{fancyhdr}
\RequirePackage{needspace}
\RequirePackage{geometry}
\RequirePackage[hidelinks]{hyperref}

\geometry{a4paper,top=17mm,bottom=18mm,left=23mm,right=23mm}
\definecolor{notegray}{gray}{0.39}
\setlength{\parindent}{0pt}
\setlength{\parskip}{2.5pt plus 1pt minus 0.5pt}
\setlist[itemize]{topsep=2pt,itemsep=1pt,parsep=0pt,leftmargin=1.45em}
\setlist[enumerate]{topsep=2pt,itemsep=2pt,parsep=0pt,leftmargin=1.65em}
\setcounter{secnumdepth}{0}
\renewcommand{\arraystretch}{1.12}

% Zahlavi: vlevo rocnik | globalni tema, vpravo cislo: nazev hodiny
% Zapati zustava zcela prazdne.
\pagestyle{fancy}
\fancyhf{}
\renewcommand{\headrulewidth}{0.35pt}
\renewcommand{\footrulewidth}{0pt}
\setlength{\headheight}{14pt}
\newcommand{\nastavPripravu}[4]{%
  \fancyhead[L]{\small #1\enspace\textbar\enspace #2}%
  \fancyhead[R]{\small #3: #4}%
}

% Minimalisticke cernobile boxy. Kazda definice zvlast.
\tcbset{commonlesson/.style={enhanced,breakable=false,colback=white,colframe=black,
  colbacktitle=white,coltitle=black,boxrule=0.45pt,
  arc=0.8mm,outer arc=0.8mm,left=3mm,right=3mm,top=1.6mm,bottom=1.6mm,
  toptitle=1mm,bottomtitle=0.7mm,fonttitle=\bfseries,
  attach boxed title to top left={xshift=2.5mm,yshift=-1.15mm},
  boxed title style={colback=white,colframe=white,boxrule=0pt,left=1mm,right=1mm}}}
\newtcolorbox{zadaniBox}[1][]{commonlesson,title={Zadání},#1}
\newtcolorbox{prikladBox}[1][]{commonlesson,title={Příklad},#1}
\newtcolorbox{definiceBox}[1][]{commonlesson,title={Definice},#1}

% Spolecny box pro zadani a jeho reseni, bez mezery mezi castmi.
% Pouziti:
%   \begin{prikladReseniBox}
%     Zadani ulohy.
%     \tcblower
%     \begin{align*} ... \end{align*}
%   \end{prikladReseniBox}
% Spodni cast ma jemny sedy podklad a delici linku; lze v ni sazet
% vice radku matematiky. Obe casti maji jen jeden spolecny ramecek.
\tcbset{solutionlesson/.style={
  commonlesson,
  enhanced jigsaw,
  breakable,
  colbacklower=black!3,
  segmentation style={solid,draw=black!25,line width=0.35pt},
  fontlower=\small\color{notegray},
  before lower={\noindent\textbf{Řešení:}\enspace},
  leftlower=3mm,rightlower=3mm
}}
\newtcolorbox{prikladReseniBox}[1][]{solutionlesson,title={Příklad},#1}
\newtcolorbox{zadaniReseniBox}[1][]{solutionlesson,title={Zadání},#1}

% Pomucky pro autora pripravy: sedive poznamky; \reseni pro starsi pripravy
\newcommand{\poznamka}[1]{\par\begingroup\small\color{notegray}\itshape
  \textnormal{\textbf{Poznámka pro učitele.}}\enspace #1\par\endgroup}
\newcommand{\reseni}[1]{\par\begingroup\small\color{notegray}
  \textbf{Očekávané řešení:}\enspace #1\par\endgroup}
\newcommand{\blok}[1]{\needspace{5\baselineskip}\subsection*{#1}}
\newcommand{\vyrok}[1]{\uv{#1}}

\usepackage{tikz}
\nastavPripravu{1. ročník}{Funkce}{67}{Lineární funkce v praxi}
\begin{document}
\blok{Od předpisu k modelu}
V hodinách 63–66 jsme počítali s předpisem $y=ax+b$ a jeho grafem. V praxi musí mít argument i hodnota jednotku a obor odpovídá smyslu veličiny. Koeficient $a$ vyjadřuje změnu výstupu na jednu jednotku vstupu; koeficient $b$ je hodnota při nulovém vstupu, pokud nula do modelu patří. Model pracuje jen v rozsahu, v němž se předpoklady uplatní.
\begin{definiceBox}
\textbf{Lineární model} veličiny $y$ podle veličiny $x$ má tvar $y=ax+b$ na určeném věcném oboru. Koeficient $a$ má jednotku „jednotka $y$ na jednotku $x$“, koeficient $b$ má stejnou jednotku jako $y$.
\end{definiceBox}
\blok{Výrobní náklady}
\begin{prikladReseniBox}
Dílna má fixní náklady $600$ Kč a za každý vyrobený kus platí dalších $75$ Kč. Zapište náklady $C(n)$, určete smysluplný obor, vypočtěte $C(8)$ a interpretujte výsledek.
\tcblower
Počet kusů $n$ je celé nezáporné číslo, tedy $n\in\mathbb N_0$. Model je $C(n)=75n+600$ Kč. Číslo $75$ má jednotku Kč/kus, $600$ Kč jsou náklady i při $n=0$. Pro osm kusů $C(8)=75\cdot8+600=1200$ Kč. Předpokládáme stejné náklady na každý kus; při velkých objemech výroby se model může přestat hodit.
\end{prikladReseniBox}
\begin{prikladReseniBox}
Kdy v předchozím modelu náklady poprvé dosáhnou alespoň $1500$ Kč? Rozlišujte spojitý výpočet a celé kusy.
\tcblower
Nejprve zjistíme, kdy náklady dosáhnou právě $1500$ Kč: $75n+600=1500$, tedy $75n=900$ a $n=12$. Náklady s každým dalším kusem rostou o $75$ Kč, proto jde zároveň o první počet kusů, při němž dosáhnou alespoň $1500$ Kč. Kontrola: $C(11)=1425$ Kč, $C(12)=1500$ Kč. Závěr „po 12 kusech“ má smysl; desetinný počet kusů by zde přípustný nebyl.
\end{prikladReseniBox}
\blok{Rovnoměrný pohyb}
\begin{prikladReseniBox}
Vlak je na začátku měření $18$ km od výchozí stanice a pokračuje stálou rychlostí $60$ km/h. Sestavte model vzdálenosti $s(t)$ po $t$ hodinách pro $0\le t\le2$. Určete $s(0{,}5)$ a význam parametrů.
\tcblower
$s(t)=60t+18$ km pro $t\in\langle0;2\rangle$. Směrnice $60$ km/h je rychlost, úsek $18$ km je počáteční vzdálenost. Při $t=0{,}5$ h vyjde $s(0{,}5)=60\cdot0{,}5+18=48$ km. Model předpokládá stálou rychlost a platí jen pro uvedené dvě hodiny; bez další informace jej neprodlužujeme do libovolné budoucnosti.
\end{prikladReseniBox}
\begin{center}
\begin{tikzpicture}[x=2.0cm,y=.035cm,>=stealth]
\draw[->] (0,0)--(2.45,0) node[right] {$t$ [h]};\draw[->] (0,0)--(0,155) node[above] {$s$ [km]};
\foreach \x in {1,2}{\draw (\x,2)--(\x,-2) node[below] {\small $\x$};}
\foreach \y in {18,48,78,138}{\draw (.03,\y)--(-.03,\y) node[left] {\small $\y$};}
\draw[very thick] (0,18)--(2,138);\fill (0,18) circle (2pt);\fill (2,138) circle (2pt);
\fill (0.5,48) circle (2pt);\node[right] at (.5,48) {\small $(0{,}5;48)$};
\end{tikzpicture}
\end{center}
Graf má vodorovnou osu v hodinách a svislou v kilometrech. Krajní body $(0;18)$ a $(2;138)$ patří do modelu. Přímku mimo časový interval nekreslíme jako tvrzení o skutečné jízdě.
\poznamka{U praktických úloh nevynechávej jednotky a věcný obor. Počet výrobků je diskrétní, čas obvykle spojitý. Výsledek dosazení mimo obor může být početně správný, ale nevypovídá o modelované situaci.}
\blok{Samostatné procvičení}
\begin{zadaniReseniBox}
\begin{enumerate}[label=\alph*)]
\item Taxi účtuje nástupní částku $90$ Kč a $25$ Kč za kilometr. Zapište cenu $T(d)$ pro $0\le d\le20$ km a určete cenu za $6$ km.
\item Nádrž obsahuje $120$ litrů vody a odtéká $8$ litrů za minutu. Sestavte objem $V(t)$ a určete, kdy je prázdná.
\item U výrobního modelu $C(n)=75n+600$ vysvětlete význam rozdílu $C(n+1)-C(n)$.
\end{enumerate}
\tcblower
\begin{enumerate}[label=\alph*)]
\item $T(d)=25d+90$ Kč, $d\in\langle0;20\rangle$ km. $T(6)=25\cdot6+90=240$ Kč. Směrnice má jednotku Kč/km.
\item $V(t)=120-8t$ litrů. Z $120-8t=0$ dostaneme $t=15$ min, proto věcný obor je $\langle0;15\rangle$. Kontrola $V(15)=0$; záporné objemy nepřipouštíme.
\item $C(n+1)-C(n)=75(n+1)+600-(75n+600)=75$ Kč: každý další kus zvýší náklady o $75$ Kč.
\end{enumerate}
\end{zadaniReseniBox}
\blok{Volitelná rezerva}
\begin{zadaniReseniBox}
Teplota se během měření od $t=0$ do $t=4$ hodin rovnoměrně zvýšila z $12\,^\circ\mathrm C$ na $20\,^\circ\mathrm C$. Sestavte model a určete teplotu po $1{,}5$ hodině.
\tcblower
Změna je $8\,^\circ\mathrm C$ za $4$ h, tedy $a=2\,^\circ\mathrm C/\mathrm h$. Model je $T(t)=12+2t$ pro $t\in\langle0;4\rangle$. $T(1{,}5)=12+3=15\,^\circ\mathrm C$. Kontrola $T(4)=20\,^\circ\mathrm C$.
\end{zadaniReseniBox}
\blok{Domácí práce}
\begin{zadaniReseniBox}
\begin{enumerate}[label=\alph*)]
\item Půjčovna účtuje $200$ Kč pevně a $50$ Kč za hodinu. Zapište cenu $P(t)$ pro celé hodiny $t=0,1,\ldots,8$ a vypočtěte $P(5)$.
\item Cyklista je $4$ km od startu a jede stálou rychlostí $12$ km/h další $3$ hodiny. Určete $s(t)$ a $s(2)$.
\item U modelu $V(t)=50-5t$ litrů, kde $t\ge0$ je čas v minutách, stanovte interval, v němž objem zůstává nezáporný.
\end{enumerate}
\tcblower
\begin{enumerate}[label=\alph*)]
\item $P(t)=200+50t$ Kč na uvedeném diskrétním oboru; $P(5)=450$ Kč.
\item $s(t)=4+12t$ km, $t\in\langle0;3\rangle$ h; $s(2)=28$ km. Směrnice má jednotku km/h.
\item $50-5t\ge0$ dává $t\le10$; spolu s $t\ge0$ je věcný interval $\langle0;10\rangle$ minut. Na konci je objem $0$.
\end{enumerate}
\end{zadaniReseniBox}
\poznamka{Jádro 45 min: význam koeficientů a jednotek 6, výrobní model 10, pohyb a graf 11, pokus žáků 9, společná kontrola 6, organizace 3 min. Rezerva 5–8 min, domácí 10–20 min. Hodina 68 porovná dvě lineární závislosti v jedné soustavě.}
\end{document}
